% =====================================================================
%  Выводы по лабораторной работе.
% =====================================================================

\labpart{Выводы}

\begin{quote}
\textit{Раздел будет окончательно сформулирован после анализа данных
из раздела «Анализ запросов из сети на SSH сервер и веб-сервер» --
семантическая и прагматическая составляющие ниже приведены в виде
каркаса и будут дополнены конкретными числами и наблюдениями.}
\end{quote}

Семантически лабораторная работа должна подтвердить (или опровергнуть)
теоретическое утверждение задания о том, что любой сервер с публичным
IP-адресом подвергается непрерывному фоновому сканированию и подбору
паролей вне зависимости от того, насколько он приметен или содержит ли
что-либо ценное для атакующего -- сам факт присутствия в интернете уже
является достаточным поводом. \textit{[после сбора данных: подтвердить
конкретными цифрами -- через сколько времени после публикации домена
появились первые обращения, сколько уникальных источников набралось за
период наблюдения и т.\,д.]}

С прагматической точки зрения работа должна показать, что настройка
журналирования -- не разовое действие, а условие, без которого
дальнейший анализ и, тем более, лабораторная работа №4.2 (сравнение
«до» и «после» установки IDS/IPS) невозможны в принципе: без
включённого заранее логирования сравнивать не с чем. \textit{[после
сбора данных: указать, какие именно практические решения по
результатам анализа логов признаны наиболее эффективными и почему]}
